\section{Introduction}
\label{sec:intro}

Automated surface inspection is the first line of defence in high-precision
manufacturing (solar cells, silicon wafers, turbine blades) and in structural
health monitoring of concrete. Field imagery is degraded by shot noise, motion
blur and non-uniform illumination, so a robust inspection pipeline must combine
noise attenuation, contrast enhancement and sub-pixel-stable edge localisation.

This work builds that pipeline from scratch, following the classical
formulations of Canny~\cite{canny1986} and Gaussian scale
space~\cite{witkin1983,lindeberg1994}. Every operator is derived from its
defining equation in pure vectorised NumPy~\cite{harris2020numpy}; OpenCV is
restricted to image \emph{I/O}. The contributions are:
(i) strided \code{conv2d}/\code{conv3d} engines that satisfy the discrete
spatial-size formula exactly, as a single \code{einsum} over a sliding-window
view (\cref{sec:conv}); (ii) a closed-form unit-gain Gaussian, unsharp masking,
and max/average pooling with analytical back-propagation matrices
(\cref{sec:enh}); (iii) a four-stage Canny detector --- normalised Sobel,
vectorised 4-sector NMS, and 8-connected hysteresis (two implementations,
asserted equal) --- in \cref{sec:canny}; and (iv) stage-wise ablations with
tolerance-aware $P/R/F_1$ against ground truth, a $(\Thi,\Tlo)$ sensitivity
grid, a check of scale-space causality, and two failure-mode analyses
(\cref{sec:results,sec:failure}).
